%============================================================
% CH2 · P2 · PROBLEM 3 — KI (AtCoder ABC138 D)
%============================================================
\subsection{Problem 2.2.3: Ki}
\problemlabel{2.2.3}{ssec:ch2-ki}

% sample 1 tree, per-vertex styles g2p2-k<k>
\tikzset{
    g2p2-k1/.style=gv, g2p2-k2/.style=gv, g2p2-k3/.style=gv, g2p2-k4/.style=gv,
    g2p2-K/.pic={
        \node[g2p2-k1] (k1) at (0,0) {1};
        \node[g2p2-k2] (k2) at (0,-1.1) {2};
        \node[g2p2-k3] (k3) at (-0.9,-2.2) {3}; \node[g2p2-k4] (k4) at (0.9,-2.2) {4};
        \foreach \a/\b in {1/2, 2/3, 2/4} \draw[gtree] (k\a) -- (k\b);
    }
}

\problemcard{2.2.3}{Ki}{AtCoder ABC138 D}{https://atcoder.jp/contests/abc138/tasks/abc138_d}{\exeasy\ (AtCoder ABC D, 400 points; the push-down direction)}{Record each operation only at the top vertex $p$; one pre-order pass pushes the totals down: $\mathrm{cnt}[v] \mathrel{+}= \mathrm{cnt}[\mathrm{par}(v)]$; iterative because $N = 2\cdot10^5$}

\subsubsection*{1. Problem statement}

A tree with $N$ vertices is rooted at vertex 1; the $i$-th edge joins $a_i$ and $b_i$. Every vertex has a counter, initially 0. Operation $j$ adds $x_j$ to the counter of \textbf{every vertex in the subtree of $p_j$}. Print all counters after the $Q$ operations.

\textbf{Constraints:} $2 \le N \le 2\cdot10^5$, $1 \le Q \le 2\cdot10^5$, $1 \le a_i < b_i \le N$, $1 \le p_j \le N$, $1 \le x_j \le 10^4$.

\begin{table}[H]
\centering\small
\begin{tabular}{@{}lll@{}}
\toprule
\textbf{Input} & \textbf{Official output} & \textbf{Our output} \\
\midrule
\texttt{4 3 / 1 2 / 2 3 / 2 4 / 2 10 / 1 100 / 3 1} & \texttt{100 110 111 110} & \texttt{100 110 111 110} \\
\texttt{6 2 / 1 2 / 1 3 / 2 4 / 3 6 / 2 5 / 1 10 / 1 10} & \texttt{20 20 20 20 20 20} & \texttt{20 20 20 20 20 20} \\
\bottomrule
\end{tabular}
\end{table}

\begin{diagrambox}
\begin{tikzpicture}[x=1cm, y=0.95cm]
    \begin{scope}
        \pic {g2p2-K};
        \node[font=\scriptsize\bfseries] at (0,0.85) {sample 1};
        \node[note, anchor=west, text width=3.9cm, align=left] at (1.4,-1.1) {op 1: $+10$ on $\{2,3,4\}$\\op 2: $+100$ on all\\op 3: $+1$ on $\{3\}$};
    \end{scope}
    \begin{scope}[xshift=6.4cm]
        \tikzset{g2p2-k1/.style=gvdone, g2p2-k2/.style=gvdone, g2p2-k3/.style=gvdst, g2p2-k4/.style=gvdone}
        \pic {g2p2-K};
        \foreach \v/\c in {1/100, 2/110, 3/111, 4/110} \node[gdist, text=deepblue, right=1pt of k\v] {\c};
        \node[font=\scriptsize\bfseries] at (0,0.85) {final counters};
    \end{scope}
    \begin{scope}[xshift=11.2cm, y=0.95cm]
        \node[gv] (m1) at (0,0) {1};
        \node[gv] (m2) at (-0.8,-1.1) {2}; \node[gv] (m3) at (0.8,-1.1) {3};
        \node[gv] (m4) at (-1.3,-2.2) {4}; \node[gv] (m5) at (-0.3,-2.2) {5}; \node[gv] (m6) at (0.8,-2.2) {6};
        \foreach \a/\b in {1/2, 1/3, 2/4, 2/5, 3/6} \draw[gtree] (m\a) -- (m\b);
        \node[font=\scriptsize\bfseries] at (0,0.85) {sample 2: $2 \times (+10)$ at 1};
    \end{scope}
\end{tikzpicture}
\figcap{The samples. In sample 2 both operations hit the root, so every counter ends at 20.}
\end{diagrambox}

\subsubsection*{2. The graph model}

\begin{modelbox}[Graph Model --- Ki]
\renewcommand{\arraystretch}{1.35}
\begin{tabularx}{\linewidth}{@{}>{\bfseries}l X@{}}
Vertices & the $N$ vertices, root 1; $\mathrm{add}[p] = \sum_{j : p_j = p} x_j$, the total put on $p$ itself \\
Edges & $N - 1$ undirected edges; $a_i < b_i$ is only a normalisation of the pair, \emph{not} ``$a_i$ is the parent'' \\
Weights & the vertex weights $\mathrm{add}[v]$ \\
Question & $\mathrm{cnt}[v] = \sum_{u \text{ ancestor of } v \text{ (or } v)} \mathrm{add}[u]$ --- a sum along the \textbf{root path}, i.e.\ push-down: $\mathrm{cnt}[v] = \mathrm{add}[v] + \mathrm{cnt}[\mathrm{par}(v)]$ \\
Size & $N, Q \le 2\cdot10^5$: \bigO{N + Q}; a path is legal, so the depth is up to $2\cdot10^5$; counters up to $Q \cdot \max x = 2\cdot10^9$ \\
\end{tabularx}
\end{modelbox}

\begin{diagrambox}
\begin{tikzpicture}[x=1cm, y=0.95cm]
    \begin{scope}
        \tikzset{g2p2-k1/.style=gvsel, g2p2-k2/.style=gvsel, g2p2-k3/.style=gvsel}
        \pic {g2p2-K};
        \node[gdist, text=softpurple, left=1pt of k1] {$+100$};
        \node[gdist, text=softpurple, left=1pt of k2] {$+10$};
        \node[gdist, text=softpurple, left=1pt of k3] {$+1$};
        \node[font=\scriptsize\bfseries] at (0,0.85) {step 1: mark only the tops};
    \end{scope}
    \draw[flow] (2.0,-1.1) -- (3.6,-1.1);
    \begin{scope}[xshift=5.6cm]
        \pic {g2p2-K};
        \draw[->, thick, accentorange] (k1) to[bend left=35] node[lbl, right, text=accentorange]{$+100$} (k2);
        \draw[->, thick, accentorange] (k2) to[bend right=30] node[lbl, left, text=accentorange]{$+110$} (k3);
        \draw[->, thick, accentorange] (k2) to[bend left=30] node[lbl, right, text=accentorange]{$+110$} (k4);
        \node[font=\scriptsize\bfseries] at (0,0.85) {step 2: push down once};
    \end{scope}
    \node[note, anchor=west, text width=4.2cm, align=left] at (8.3,-1.1) {a vertex $v$ is in the subtree of $p$ $\iff$ $p$ is on the root path of $v$; so $v$ receives exactly the marks on its root path};
\end{tikzpicture}
\figcap{The model: ``add to a subtree'' turned around becomes ``sum over the root path'', which is what a push-down computes.}
\end{diagrambox}

\subsubsection*{3. Key observations}

\begin{itemize}[leftmargin=*]
    \item \textbf{Turn the operation around.} $v \in \mathrm{subtree}(p) \iff p \text{ is an ancestor of } v$ (or $v$ itself). So the final counter of $v$ is the sum of the marks $\mathrm{add}[u]$ over its root path --- the pre-order direction of the pattern.
    \item \textbf{Operations commute.} All operations are additions and nothing is asked in between, so we may postpone all the work: record \bigO{1} per operation, then one pass.
    \item \textbf{Root it yourself.} $a_i < b_i$ does not say who is the parent: with edges $1{-}3$, $2{-}3$ vertex 3 is the parent of 2. A program that sets $\mathrm{par}[b_i] = a_i$ prints \texttt{0 0 1} instead of \texttt{0 1 1} for the operation ``$+1$ on 3''.
    \item \textbf{Depth.} A path of $2\cdot10^5$ vertices is a legal tree; our recursive push-down died on it without output. BFS order makes the pass iterative.
\end{itemize}

\subsubsection*{4. Approaches}

\begin{bruteforcebox}
\textbf{Idea: do what the statement says.} Root the tree once, then for each operation walk the subtree of $p_j$ (downwards only) and add $x_j$.

\begin{lstlisting}[style=cpp]
#include <bits/stdc++.h>
using namespace std;
typedef long long ll;

// Brute force: perform every operation literally. Root the tree once (parent array), then for each
// operation walk the whole subtree of p (only downwards: never back to the parent) and add x.
int main() {
    int n, q;
    cin >> n >> q;
    vector<vector<int>> adj(n + 1);
    for (int i = 0; i < n - 1; i++) {
        int a, b;
        cin >> a >> b;
        adj[a].push_back(b);
        adj[b].push_back(a);
    }
    vector<int> parent(n + 1, 0), order(1, 1);
    for (size_t i = 0; i < order.size(); i++)
        for (int v : adj[order[i]])
            if (v != parent[order[i]]) { parent[v] = order[i]; order.push_back(v); }
    vector<ll> cnt(n + 1, 0);
    for (int j = 0; j < q; j++) {
        int p, x;
        cin >> p >> x;
        vector<int> st(1, p);
        while (!st.empty()) {
            int u = st.back(); st.pop_back();
            cnt[u] += x;
            for (int v : adj[u]) if (v != parent[u]) st.push_back(v);
        }
    }
    for (int v = 1; v <= n; v++) cout << cnt[v] << (v == n ? '\n' : ' ');
}
\end{lstlisting}

\begin{complexitybox}
\textbf{Time:} \bigO{\sum_j \mathrm{sz}[p_j]} $= \bigO{NQ}$ $= 4\cdot10^{10}$ when every operation hits the root. On a path of $2\cdot10^5$ vertices with $2\cdot10^5$ operations at the root it had not finished after three minutes. \textbf{Space:} \bigO{N}. \textbf{Verdict:} \tle; reference for testing.
\end{complexitybox}
\end{bruteforcebox}

\begin{algorithmbox}[Approach 2: mark the top, push down in BFS order]
\textbf{Idea.} $\mathrm{cnt}[p_j] \mathrel{+}= x_j$ for each operation. Then a BFS from 1 that skips the parent: when it discovers $v$ from $u$, $\mathrm{cnt}[u]$ is final, and $\mathrm{cnt}[v] \mathrel{+}= \mathrm{cnt}[u]$. This is the forward loop of template (2), fused with the BFS itself.

\begin{lstlisting}[style=cpp]
#include <bits/stdc++.h>
using namespace std;
typedef long long ll;

// Ki (ABC138 D): put x on the counter of p only, then push the counters DOWN the tree:
// final[v] = add[v] + final[parent(v)], in an order where every parent comes first (BFS order).
int main() {
    ios::sync_with_stdio(false);
    cin.tie(nullptr);
    int n, q;
    cin >> n >> q;
    vector<vector<int>> adj(n + 1);
    for (int i = 0; i < n - 1; i++) {
        int a, b;
        cin >> a >> b;
        adj[a].push_back(b);
        adj[b].push_back(a);
    }
    vector<ll> cnt(n + 1, 0);
    for (int j = 0; j < q; j++) {
        int p, x;
        cin >> p >> x;
        cnt[p] += x;                               // only the top of the subtree, O(1)
    }
    vector<int> parent(n + 1, 0), order(1, 1);
    parent[1] = -1;
    for (size_t i = 0; i < order.size(); i++) {    // BFS from the root, no recursion
        int u = order[i];
        for (int v : adj[u])
            if (v != parent[u]) {
                parent[v] = u;
                cnt[v] += cnt[u];                  // pre-order push-down: u is already final
                order.push_back(v);
            }
    }
    for (int v = 1; v <= n; v++) cout << cnt[v] << (v == n ? '\n' : ' ');
}
\end{lstlisting}

\begin{complexitybox}
\textbf{Time:} \bigO{N + Q}. \textbf{Space:} \bigO{N}. \textbf{Verdict:} reproduces both samples; on the path of $2\cdot10^5$ vertices with all operations at the root it prints $2\cdot10^9$ everywhere in $0.9$\,s wall time here, start-up included.
\end{complexitybox}
\end{algorithmbox}

\begin{invariantbox}[Why one push-down is enough]
\textbf{Claim.} When $v$ is discovered from $u = \mathrm{par}(v)$, $\mathrm{cnt}[u] = \sum_{w \in \mathrm{path}(1, u)} \mathrm{add}[w]$, and after the update $\mathrm{cnt}[v] = \sum_{w \in \mathrm{path}(1, v)} \mathrm{add}[w]$.

\emph{Induction along the BFS order.} The root is never updated: $\mathrm{cnt}[1] = \mathrm{add}[1]$, its root path is $\{1\}$. A vertex $u$ is taken from the queue only after it was discovered (so its own update is done), and $\mathrm{cnt}[u]$ never changes afterwards: only its parent writes into it, once. So $\mathrm{cnt}[v] = \mathrm{add}[v] + \mathrm{cnt}[u]$ is the sum over $\mathrm{path}(1, u) \cup \{v\} = \mathrm{path}(1, v)$.

\emph{The answer.} $\mathrm{path}(1, v)$ is the set of $p$ whose subtree contains $v$, so $\mathrm{cnt}[v] = \sum_{j : v \in \mathrm{subtree}(p_j)} x_j$, the required counter. \hfill$\blacksquare$
\end{invariantbox}

\begin{alternativebox}[Approach 3: a difference array on pre-order intervals]
Number the vertices in DFS pre-order: $\mathrm{pos}[v]$. The subtree of $p$ is then the block of positions $\mathrm{pos}[p], \dots, \mathrm{pos}[p] + \mathrm{sz}[p] - 1$ (this is the parenthesis theorem of Lesson~\ref{ssec:ch2-timestamps}, the core of Pattern~3). ``Add $x$ to a block'' on an array is a difference update: $+x$ at the start, $-x$ just after the end; one prefix sum at the end gives every value. This version is the one that generalises when queries are mixed with updates: replace the difference array by a Fenwick tree (Chapter~15).

\begin{lstlisting}[style=cpp]
#include <bits/stdc++.h>
using namespace std;
typedef long long ll;

// Alternative: number the vertices in DFS pre-order. The subtree of p is the block of positions
// pos[p] .. pos[p] + sz[p] - 1, so "add x to the subtree" becomes a range update on an array:
// diff[pos[p]] += x, diff[pos[p] + sz[p]] -= x, and one prefix sum at the end.
int main() {
    ios::sync_with_stdio(false);
    cin.tie(nullptr);
    int n, q;
    cin >> n >> q;
    vector<vector<int>> adj(n + 1);
    for (int i = 0; i < n - 1; i++) {
        int a, b;
        cin >> a >> b;
        adj[a].push_back(b);
        adj[b].push_back(a);
    }
    // iterative DFS: pre-order positions, then sizes from the reversed pre-order
    vector<int> pos(n + 1), parent(n + 1, 0), pre, st(1, 1);
    while (!st.empty()) {
        int u = st.back(); st.pop_back();
        pos[u] = pre.size();
        pre.push_back(u);
        for (int v : adj[u])
            if (v != parent[u]) { parent[v] = u; st.push_back(v); }
    }
    vector<int> sz(n + 1, 1);
    for (int i = n - 1; i >= 1; i--) sz[parent[pre[i]]] += sz[pre[i]];
    vector<ll> diff(n + 1, 0);
    for (int j = 0; j < q; j++) {
        int p, x;
        cin >> p >> x;
        diff[pos[p]] += x;                         // the block of p starts here ...
        diff[pos[p] + sz[p]] -= x;                 // ... and ends just before here
    }
    for (int i = 1; i < n; i++) diff[i] += diff[i - 1];
    for (int v = 1; v <= n; v++) cout << diff[pos[v]] << (v == n ? '\n' : ' ');
}
\end{lstlisting}
The stack here pushes all children of the popped vertex at once. In a general graph that is the ``shortcut that is not DFS'' of Lesson~\ref{ssec:ch2-call-stack}; in a tree every vertex is pushed exactly once (by its parent), and the last-pushed child's whole subtree is popped before its siblings, so the subtrees still occupy contiguous blocks --- a genuine pre-order, with the children in reverse order.
\begin{complexitybox}
\textbf{Time:} \bigO{N + Q}. \textbf{Space:} \bigO{N}. \textbf{Verdict:} reproduces both samples and prints the same output as Approach~2 on the large tests (identical checksums).
\end{complexitybox}
\end{alternativebox}

\begin{tricksbox}[Approach 4: Implementation Refinements]
\begin{itemize}[leftmargin=*]
    \item \textbf{\texttt{int} or \texttt{long long}?} The largest counter is $Q \cdot \max x = 2\cdot10^5 \cdot 10^4 = 2\cdot10^9 < 2^{31} - 1 = 2\,147\,483\,647$: \texttt{int} fits, with 7\,\% to spare. One more zero in any limit breaks it; \texttt{long long} costs nothing here.
    \item \textbf{Fast I/O.} $8\cdot10^5$ numbers in, $2\cdot10^5$ out: \texttt{sync\_with\_stdio(false)} and one \texttt{'\textbackslash n'} at the end.
    \item \textbf{Fusing passes.} The push-down can be done inside the BFS (as above) because a vertex is final at the moment it is dequeued. Aggregation (the up direction) cannot be fused this way: it needs the backward loop after the BFS.
\end{itemize}
\end{tricksbox}

\subsubsection*{5. Trace}

\begin{tracebox}[{Sample 1: edges $1{-}2$, $2{-}3$, $2{-}4$; operations $(2, 10)$, $(1, 100)$, $(3, 1)$}]
\begin{center}
\begin{tikzpicture}[x=0.95cm, y=0.85cm]
    \tikzset{gv/.append style={minimum size=6.5mm, font=\scriptsize\bfseries}}
    \begin{scope}
        \pic {g2p2-K};
        \foreach \v/\c in {1/100, 2/10, 3/1, 4/0} \node[gdist, text=softpurple, right=1pt of k\v] {\c};
        \node[font=\scriptsize\bfseries] at (0,0.8) {1. marks after reading};
    \end{scope}
    \begin{scope}[xshift=4.2cm]
        \tikzset{g2p2-k1/.style=gvdone, g2p2-k2/.style=gvcur}
        \pic {g2p2-K};
        \foreach \v/\c in {1/100, 2/110, 3/1, 4/0} \node[gdist, text=deepblue, right=1pt of k\v] {\c};
        \node[font=\scriptsize\bfseries] at (0,0.8) {2. from 1: $\mathrm{cnt}[2] = 10 + 100$};
    \end{scope}
    \begin{scope}[xshift=8.4cm]
        \tikzset{g2p2-k1/.style=gvdone, g2p2-k2/.style=gvdone, g2p2-k3/.style=gvcur, g2p2-k4/.style=gvcur}
        \pic {g2p2-K};
        \foreach \v/\c in {1/100, 2/110, 3/111, 4/110} \node[gdist, text=deepblue, right=1pt of k\v] {\c};
        \node[font=\scriptsize\bfseries] at (0,0.8) {3. from 2: 3 and 4};
    \end{scope}
\end{tikzpicture}
\end{center}
\begin{center}\small
\begin{tabular}{@{}clll@{}}
\toprule
\textbf{dequeue $u$} & \textbf{discover} & \textbf{update} & \textbf{$\mathrm{cnt}[1..4]$ after} \\
\midrule
--- & --- & marks: $\mathrm{cnt}[2] = 10$, $\mathrm{cnt}[1] = 100$, $\mathrm{cnt}[3] = 1$ & $100, 10, 1, 0$ \\
1 & 2 & $\mathrm{cnt}[2] = 10 + 100$ & $100, 110, 1, 0$ \\
2 & 3, 4 & $\mathrm{cnt}[3] = 1 + 110$, $\mathrm{cnt}[4] = 0 + 110$ & $100, 110, 111, 110$ \\
3, 4 & --- & leaves & $100, 110, 111, 110$ \\
\bottomrule
\end{tabular}
\end{center}
\textbf{Approach 3 on the same sample.} The stack DFS gives the pre-order $1, 2, 4, 3$ (children of 2 pushed as 3, 4; 4 is popped first), so $\mathrm{pos} = (0, 1, 3, 2)$ for vertices $1..4$ and $\mathrm{sz} = (4, 3, 1, 1)$.
\begin{center}
\begin{tikzpicture}[x=1.15cm]
    \node[font=\scriptsize\bfseries, anchor=east] at (-0.6,0.75) {vertex};
    \foreach \v [count=\i from 0] in {1,2,4,3} \node[idx] at (\i,0.75) {\v};
    \node[font=\scriptsize\bfseries, anchor=east] at (-0.6,0) {diff};
    \foreach \d [count=\i from 0] in {100,10,0,1,$-111$} \node[cell, minimum width=10mm] (d\i) at (\i,0) {\d};
    \node[font=\scriptsize\bfseries, anchor=east] at (-0.6,-0.75) {prefix};
    \foreach \d [count=\i from 0] in {100,110,110,111} \node[cell, minimum width=10mm, fill=computedfill] at (\i,-0.75) {\d};
    \draw[decorate, decoration={brace, amplitude=4pt, mirror}, primaryblue] (0.6,-1.2) -- (3.4,-1.2) node[midway, below=4pt, font=\scriptsize]{subtree(2) = positions 1..3: $+10$ at 1, $-10$ at 4};
    \node[note, anchor=west, text width=3.6cm, align=left] at (5.4,0) {op (1,100): $+100$ at 0, $-100$ at 4\\op (3,1): $+1$ at 3, $-1$ at 4};
\end{tikzpicture}
\end{center}
Read at $\mathrm{pos}[v]$: vertex 1 $\to 100$, 2 $\to 110$, 3 $\to 111$ (position 3), 4 $\to 110$ (position 2). The arrays $\mathrm{pos}$, $\mathrm{sz}$, diff and prefix were printed by an instrumented copy of Approach~3.
\end{tracebox}

\subsubsection*{6. Test suite}

\begin{testbox}
\centering\small
\begin{tabular}{@{}p{5.4cm}p{3.2cm}p{4.8cm}@{}}
\toprule
\textbf{Input} & \textbf{Output} & \textbf{Why} \\
\midrule
\texttt{2 1 / 1 2 / 2 5} & \texttt{0 5} & operation at a leaf \\
\texttt{4 2 / 1 4 / 2 4 / 3 4 / 4 7 / 1 1} & \texttt{1 8 8 8} & star centred at 4, but the root 1 is a leaf \\
\texttt{5 2 / 1 2 / 2 3 / 3 4 / 4 5 / 3 10 / 5 1} & \texttt{0 0 10 10 11} & path, operations in the middle \\
\texttt{3 2 / 1 3 / 2 3 / 3 1 / 2 2} & \texttt{0 3 1} & $a_i < b_i$ but 3 is the parent of 2 \\
path $1..2\cdot10^5$, $2\cdot10^5 \times (1, 10^4)$ & $2\cdot10^9$ everywhere & deepest tree, largest counter \\
random tree, $N = Q = 2\cdot10^5$ & (same checksum) & Approaches 2 and 3 agree \\
\bottomrule
\end{tabular}
\end{testbox}
All outputs from Approach~2; Approaches 1 and 3 printed the same on the small tests. In addition, the three approaches agreed on $3000$ random inputs ($N \le 10$, $Q \le 8$, random / path / star / caterpillar trees with shuffled labels and randomly ordered edges).

\subsubsection*{7. Common mistakes}

\begin{warningbox}[Mistake 1 --- trusting $a_i < b_i$]
\begin{lstlisting}[style=snippet]
for (each edge a b) par[b] = a;                       // WRONG: smaller label != parent
for (int v = 2; v <= n; v++) cnt[v] += cnt[par[v]];
\end{lstlisting}
Both samples happen to have $a_i = \mathrm{par}(b_i)$, so this passes them. On \texttt{3 1 / 1 3 / 2 3 / 3 1} it prints \texttt{0 0 1} instead of \texttt{0 1 1}; our random test found a mismatch on its very first input.
\end{warningbox}

\begin{warningbox}[Mistake 2 --- recursion on a path]
A recursive \texttt{dfs(v, p)} doing the same push-down crashed (no output at all) on the path of $2\cdot10^5$ vertices with the default 2\,MB stack; it works on random trees, whose depth is small. Use the BFS order.
\end{warningbox}

\begin{warningbox}[Mistake 3 --- simulating the operations]
Walking the subtree per operation looks harmless on the samples but is \bigO{NQ} (Approach~1). The point of the problem is to defer the work.
\end{warningbox}

\subsubsection*{8. Takeaways}

\begin{itemize}[leftmargin=*]
    \item ``Add to the subtree of $p$'' $=$ ``every descendant of $p$ receives it'' $=$ ``$v$ receives the marks on its root path'': mark the top, push down once.
    \item Push-down is pre-order: forward over a BFS order, and it can be fused into the BFS.
    \item Never derive the parent from the input order of an edge; root the tree yourself.
    \item Subtrees are contiguous in pre-order: subtree updates become range updates (Pattern~3, Chapter~15).
\end{itemize}

\begin{variantsbox}
\variant{Queries between updates}{\varnone}
  {updates ``add $x$ to subtree($p$)'' and queries ``value of $v$'' are interleaved.}
  {the same pre-order blocks.}
  {a Fenwick tree over positions: range add, point query, \bigO{\log N} each (Chapter~15). Our Fenwick version matched a per-operation subtree walk on $3000$ random trees with 15 interleaved operations each.}
\variant{Add along a root path}{\varnone}
  {operation: add $x$ to every vertex on the path from $p$ up to the root.}
  {the transpose: $v$ receives $x$ iff $p$ is in the subtree of $v$.}
  {mark $p$, then \emph{aggregate up} (backward loop): $\mathrm{cnt}[v] = \sum_{\mathrm{subtree}(v)} \mathrm{add}$. Matched a per-operation walk up the parents on $3000$ random trees.}
\variant{Best path ending here}{\varin{https://codeforces.com/problemset/problem/682/C}}
  {for every vertex, the largest sum of an upward path ending at it, with negative edges.}
  {push-down of a maximum, restarted when the parent's value is negative.}
  {$\mathrm{best}[v] = \max(\mathrm{best}[\mathrm{par}], 0) + c$ (Exercise E2.4).}
\variant{Counts per label}{\varin{https://leetcode.com/problems/number-of-nodes-in-the-sub-tree-with-the-same-label/}}
  {how many vertices of the subtree share $v$'s label.}
  {a pre-order position interval per subtree.}
  {a global counter read at entry and exit: the before/after trick (Exercise E2.1).}
\end{variantsbox}
